problem:  
Let \(M^n \subset \mathbb{S}^{n+p}(1)\) be a compact **minimal submanifold with parallel second fundamental form**, i.e. an orbit of an \(s\)‑representation. For such an \(M\), the space of infinitesimal minimal deformations (Jacobi fields) decomposes into irreducible \(K\)‑modules under the isotropy group \(K\) of the corresponding symmetric space \(G/K\). In many cases (for example, certain Clifford tori), some of these components do not correspond to ambient Killing fields, yielding infinitesimal non‑isometric deformations.

**Open problem:**  
Classify all compact minimal \(s\)-representation orbits \(M \subset \mathbb{S}^{n+p}(1)\) for which **Jacobi rigidity** holds, i.e.
\[
\mathcal{K}(M) = \mathcal{K}_{\mathrm{iso}}(M),
\]
and determine the structure of the extra Jacobi field spaces in the non‑rigid cases.  
Equivalently, describe completely the **representation‑theoretic decomposition** of the Jacobi operator’s kernel on the normal bundle for every family of irreducible \(s\)-representation orbits, identifying when each component integrates to genuine minimal deformations.

Why is it a "good" problem:  
This refined problem turns the false general statement into a concrete and open **classification programme** lying at the intersection of geometry and representation theory.  
- It unifies known scattered rigidity results (for example, totally geodesic spheres, certain Veronese embeddings) and extends them into an overarching question: *which symmetric minimal embeddings are infinitesimally rigid, and why?*  
- Analytically, it links the **spectrum of the Jacobi operator** with the representation theory of the isotropy group, demanding a deep understanding of how geometric curvature translates into spectral multiplicities.  
- Representation‑theoretically, it invites the development of tools to compute normal bundle Laplacians for homogeneous spaces and distinguish Killing‑type from genuinely new deformation‑type modes.  
- Conceptually, it asks for the frontier between **intrinsic symmetry and extrinsic flexibility**, a theme central to modern submanifold and homogeneous geometry.  

The problem is **well‑posed, open, and nontrivial**, with potentially profound implications for the structure theory of minimal homogeneous submanifolds and for broader links between global analysis and symmetric space geometry.